\section{The Composition of a Small Leak}\label{sec:attack:leak}

This section asks whether the composition of a leak, at a fixed number of rows, changes word-F1.
The budgets of the previous sections count whole paired sequences; a leak may instead hold a few rows from each of many prompts, or rows of one kind of word.
The leak is fixed near \vf{90} rows, the size of ten Corpus A prompts, and varies which rows it holds.
The corpus is legal128w, whose prompts average \vf{61} words (Section~\ref{sec:hub:legal}), so that \vf{90} rows fit in one or two prompts or spread over ninety.
The hybrid fit has $\alpha = 0.99$, $\tau = d$, and $B = 8{,}192$; the supervised fit reads the same rows. % src: results/legal128w/v7_2/main/mixed-oct_w5-s1-L90.json setting (alphas 0.99, tau_sqdist 0.06485, blind_batch_rows 8192)
All results of this section rest on one seed.

\subsection{Rows spread over many prompts}\label{sec:attack:leak:spread}

The first part of H6 meets its criterion.
One leak takes one row from each of \vf{90} prompts; under the hybrid fit it yields a word-F1 of \vf{0.321}, against \vf{0.178} for \vf{90} rows from two prompts (Table~\ref{tab:attack:shape}).
The difference of \vf{0.14} exceeds the \vf{0.03} of the criterion. % src: results/legal128w/v7_2/main/mixed-oct_w5-s1-L90.json (L90-random 0.3207); mixed-oct_w5b-L-prefix.json (L90-prefix 0.1782)
No interval is stored for that difference, because the two fits come from separate runs.
One whole prompt of \vf{78} rows scores \vf{0.162}. % src: mixed-oct_w5b-L.json (b1 0.1615)
For these sampled leaks, spreading the rows over prompts raises word-F1; the comparison does not separate the number of prompts from the content of the particular prompts drawn.
Most of the difference lies between one or two prompts and ten: with nine consecutive rows from each of ten prompts, the leak scores \vf{0.304}, only \vf{0.017} below the leak over \vf{90} prompts (interval \vf{0.008} to \vf{0.025}). % src: results/legal128w/v7_2/main/mixed-oct_w5-s1-L90-ext2.json seeds.1.budgets.90.axis_effect delta_h:...L90-span9 (-0.0168 [-0.0254, -0.0084])
On the \vf{4{,}500} held-out prompts the leak over \vf{90} prompts scores \vf{0.326} under the hybrid fit and \vf{0.106} under the supervised fit, a gain of \vf{0.220} (interval \vf{0.217} to \vf{0.223}). % src: results/legal128w/v7_2/main/oct_w5-summary-confirmation-A.json levels H/random/L90/all, S/random/L90; contrasts E/random/L90/all
The supervised fit stays below \vf{0.11} for every leak of about \vf{90} rows, and for the leaks spread over \vf{10} or \vf{90} prompts the gain of the matching term is about twice the supervised score.

\begin{table}[H]
\caption{At about \vf{90} rows, a leak spread over \vf{10} or \vf{90} prompts decodes at \vf{0.30} to \vf{0.32} under the hybrid fit; a leak from one or two prompts decodes at \vf{0.16} to \vf{0.18}.
Word-F1 averaged over three decoders, Pythia on legal128w, seed \vf{1}, \vf{500} selection prompts; the hub reference is \vf{0.982}.
The difference is hybrid minus supervised with a \vf{95} percent interval; the mismatched-prompt control of the hybrid fits lies between \vf{0.06} and \vf{0.13}.
The last row holds \vf{864} rows and is not matched in size to the others.}
\label{tab:attack:shape}
\centering
\small\setlength{\tabcolsep}{4pt}
\begin{tabular}{lccccc}
\toprule
Leak & Prompts & Rows & Hybrid fit & Supervised fit & Difference \\
\midrule
One row per prompt & \vf{90} & \vf{90} & \vf{0.321} & \vf{0.108} & \vf{$+0.212$} [\vf{0.202}, \vf{0.222}] \\ % src: results/legal128w/v7_2/main/mixed-oct_w5-s1-L90.json cells *L90-random (0.3207, 0.1084, +0.2123 [0.2022, 0.2218]; control 0.1006)
Nine consecutive rows per prompt & \vf{10} & \vf{90} & \vf{0.304} & \vf{0.094} & \vf{$+0.210$} [\vf{0.200}, \vf{0.220}] \\ % src: mixed-oct_w5-s1-L90-ext2.json cells *L90-span9 (0.3039, 0.0938, +0.2101 [0.1999, 0.2197]; control 0.1043)
One prompt and part of a second & \vf{2} & \vf{90} & \vf{0.178} & \vf{0.049} & \vf{$+0.130$} [\vf{0.121}, \vf{0.138}] \\ % src: mixed-oct_w5b-L-prefix.json cells *L90-prefix (0.1782, 0.0487, +0.1295 [0.1213, 0.1378]; control 0.0605)
One whole prompt & \vf{1} & \vf{78} & \vf{0.162} & \vf{0.087} & \vf{$+0.075$} [\vf{0.067}, \vf{0.082}] \\ % src: mixed-oct_w5b-L.json cells b1/* (0.1615, 0.0867, +0.0748 [0.0669, 0.0822]; control 0.0738)
\midrule
One row per prompt & \vf{864} & \vf{864} & \vf{0.628} & \vf{0.548} & \vf{$+0.080$} [\vf{0.072}, \vf{0.088}] \\ % src: mixed-oct_w5-s1-L864.json cells *L864-random (0.6281, 0.5480, +0.0801 [0.0724, 0.0876]; control 0.1266)
\bottomrule
\end{tabular}
\end{table}

With \vf{864} rows, one from each of \vf{864} prompts, the supervised fit reaches \vf{0.548} and the gain of the matching term falls to \vf{0.080}.
On the held-out prompts the two scores are \vf{0.631} and \vf{0.546}, a gain of \vf{0.085} (interval \vf{0.082} to \vf{0.087}). % src: oct_w5-summary-confirmation-A.json levels H/random/L864/all, S/random/L864; contrasts E/random/L864/all
The gain thus shrinks as the leak grows, as it does on Corpus A (Section~\ref{sec:attack:hybrid}).

Table~\ref{tab:attack:legal-examples} shows three reconstructions under the leak of one row from each of \vf{90} prompts, through the SmolLM decoder, whose median prompt scores \vf{0.55}. % src: results/legal128w/v7_2/main/mixed-v7_9/runs/b90-a0.99-sq-m0-t0.06485-r1-B8192-L90-random-s1/endpoint/decoding.jsonl, SmolLM records (median 0.5502, maximum 0.836)
Through SmolLM no selection prompt is reconstructed exactly; the highest score is \vf{0.84}.
In the three examples the errors are fluent substitutions: ``registrant'' becomes ``insurer'', and dates, exhibit numbers, section numbers, and a company name are replaced by others of the same kind.
The reconstruction itself does not mark which words match the source and which are substitutions.

\begin{table}[H]
\caption{On legal text, the hybrid fit at \vf{90} rows recovers the frame of a passage and substitutes names and numbers.
Reconstructions through the SmolLM decoder for three legal128w selection prompts, chosen at high, intermediate, and low word-F1 (seed \vf{1}, one row from each of \vf{90} prompts).
Decoded words that match the true prompt are green; the others are red and bold.}
\label{tab:attack:legal-examples}
\centering
\footnotesize
\begin{tabular}{p{0.42\textwidth}p{0.42\textwidth}c}
\toprule
True prompt & Decoded & Word-F1 \\
\midrule
disclosure controls and procedures, as of the end of the period covered by this report based on such evaluation; and (d) Disclosed in this report any change in the registrant's internal control over financial reporting that occurred during the registrant's most recent fiscal quarter (the registrant's fourth fiscal quarter in the case of an annual report) that has materially affected, or is reasonably likely to materially affect, the registrant's internal control over financial reporting; and & \miss{responsibility} \hit{controls and procedures, as of the end of the period covered by this report based on such evaluation; and} \miss{(i) Considered} \hit{in this report any change in the} \miss{insurer's} \hit{internal control over financial reporting that occurred during the} \miss{insurer's} \hit{most recent fiscal} \miss{year} \hit{(the} \miss{insurer's first} \hit{fiscal} \miss{year} \hit{in the case of} \miss{a} \hit{annual report) that has materially affected, or is reasonably likely to materially affect, the} \miss{insurer's} \hit{internal control over financial} \miss{reporting} & \vf{0.83} \\ % src: decoding.jsonl above, law-edgar-00014181, smollm (0.8322); truth results/legal128w/v7_2/main/position/dev-truth.json
\midrule
operations for the year's endedDecember31, 2010 and 2009are filed as Exhibit 99.2 herewith. (d)Exhibits ExhibitNo. Description Stock Purchase Agreement, dated as of December23, 2010.,.(1) Audited financial statements as of and for the years ended December 31, 2010 and 2009 for CMAC, Inc. & \miss{references} \hit{for the} \miss{year} \hit{of} \miss{finished} \hit{December 2009 and} \miss{2000 are} \hit{filed as} \miss{Examination 61.2 hereafter. (A) Ex. No.} \hit{Description Stock Purchase Agreement, dated as of} \miss{June 07,2009, (B) Unaudited} \hit{financial statements as} \miss{of} \hit{and for the} \miss{months} \hit{ended} \miss{December 2009} \hit{and} \miss{2000} \hit{for} \miss{Kelebank, LLC.} & \vf{0.55} \\ % src: same file, law-edgar-00009775 (0.5455)
\midrule
Lender Credit Decision 58 Section 10.8 Expenses; Indemnities 58 Section 10.9 Resignation of Administrative Agent 59 Section 10.10 Release of Collateral or Guarantors 60 ARTICLE 11 MISCELLANEOUS 60 Section 11.1 Amendments, Waivers, Etc 60 Section 11.2 Assignments and Participations; & \hit{Lender} \miss{Creditability 31} \hit{Section} \miss{2.2 Exclusions; Unearned Benefits 26} \hit{Section} \miss{2.2} \hit{Resignation of Administrative} \miss{Officer9} \hit{Section} \miss{2.9} \hit{Release of Collateral and} \miss{Guarantors9} \hit{Section} \miss{2.9 Refusals, Exceptions9} \hit{Section} \miss{2.9} \hit{Amendments,} \miss{Revocations, Exceptions9 Section 2.9} \hit{Assignments} \miss{and Subscriptions;} & \vf{0.40} \\ % src: same file, law-edgar-00007393 (0.4000)
\bottomrule
\end{tabular}
\end{table}

\subsection{Word class of the leaked rows}\label{sec:attack:leak:words}

The second part of H6 meets its criterion: no rule that restricts the leaked words exceeds a random draw of rows by \vf{0.03}.
Each rule takes one eligible row from each of \vf{90} prompts and skips a prompt without an eligible row, so the rules can differ in the prompts they draw as well as in the words they leak (Table~\ref{tab:attack:words}). % src: src/v7_2/anchors_linear.py (leak rules: prompts without an eligible row are skipped)
Stop words are the words of a fixed stop list; mid-frequency words have rank \vf{101} to \vf{1{,}000} in the word counts of the legal128w training prompts, and rare words a rank above \vf{1{,}000}. % src: oct_w5/plan.md section 3; src/v7_2/anchors_linear.py (rank bands)
The mid-frequency and rare rules also admit groups of punctuation alone whose forms fall in those ranks. % src: src/v7_2/anchors_linear.py (mid, rare against mid_word, rare_word)
The one-word rule leaks the same mid-frequency word in all \vf{90} prompts, and the distinct rule draws rows at random but leaks no word twice.

\begin{table}[H]
\caption{No word class improves on a random draw of rows; stop-word rows and one repeated word give the hybrid fit no gain.
Word-F1 averaged over three decoders, Pythia on legal128w, seed \vf{1}, \vf{500} selection prompts, \vf{90} rows from \vf{90} prompts.
The last column is the hybrid fit under the rule minus the hybrid fit under the random draw, with a \vf{95} percent interval.}
\label{tab:attack:words}
\centering
\small\setlength{\tabcolsep}{4pt}
\begin{tabular}{lccc}
\toprule
Leaked rows & Hybrid fit & Supervised fit & Hybrid fit minus random draw \\
\midrule
Random draw & \vf{0.321} & \vf{0.108} & \\ % src: results/legal128w/v7_2/main/mixed-oct_w5-s1-L90.json cells *L90-random (0.3207, 0.1084)
Distinct words & \vf{0.329} & \vf{0.109} & \vf{$+0.009$} [\vf{0.001}, \vf{0.017}] \\ % src: mixed-oct_w5-s1-L90-ext.json cells *L90-distinct (0.3295, 0.1087); axis_effect delta_h (+0.0088 [0.0007, 0.0166])
Mid-frequency words & \vf{0.242} & \vf{0.094} & \vf{$-0.078$} [\vf{$-0.087$}, \vf{$-0.069$}] \\ % src: mixed-oct_w5-s1-L90.json cells *L90-mid (0.2423, 0.0938); delta_h (-0.0784 [-0.0871, -0.0692])
Rare words & \vf{0.226} & \vf{0.072} & \vf{$-0.095$} [\vf{$-0.104$}, \vf{$-0.086$}] \\ % src: same file, *L90-rare (0.2255, 0.0717); delta_h (-0.0952 [-0.1041, -0.0860])
Stop words & \vf{0.058} & \vf{0.056} & \vf{$-0.263$} [\vf{$-0.273$}, \vf{$-0.252$}] \\ % src: same file, *L90-stop (0.0582, 0.0555); delta_h (-0.2625 [-0.2725, -0.2516])
One repeated word & \vf{0.054} & \vf{0.057} & \vf{$-0.267$} [\vf{$-0.278$}, \vf{$-0.256$}] \\ % src: same file, *L90-one_word (0.0535, 0.0568); delta_h (-0.2672 [-0.2781, -0.2563])
\bottomrule
\end{tabular}
\end{table}

The largest positive difference is \vf{0.009}, for distinct words; every class rule lowers the hybrid fit.
Before the run, the evaluation plan predicted that mid-frequency words would raise the hybrid fit by at least \vf{0.03}; on the held-out prompts they lower it by \vf{0.074} (interval \vf{0.072} to \vf{0.077}), so that prediction is contradicted. % src: oct_w5/plan.md section 6 (P2); results/legal128w/v7_2/main/oct_w5-summary-confirmation-B.json decisions.P2 (-0.0743 [-0.0773, -0.0715], contradicted)
A leak of stop words gives the matching term nothing to add at \vf{90} rows: on the held-out prompts the hybrid and supervised fits both score \vf{0.056}. % src: oct_w5-summary-confirmation-A.json levels H/stop/L90/all (0.0564), S/stop/L90 (0.0561); contrasts E/stop/L90/all (+0.0004 [-0.0011, 0.0019])
At \vf{864} stop-word rows the hybrid fit reaches \vf{0.461} against \vf{0.198} for the supervised fit on the held-out prompts. % src: same file, levels H/stop/L864/all (0.4610), S/stop/L864 (0.1980)
The one-word rule leaks a single word and scores as low as the stop-word rule; whether the stop-word leak fails for the same reason, few distinct words, was not measured.

The word-class rules therefore lower word-F1 by as much as \vf{0.27}, more than the spread over prompts changes it (\vf{0.14}), and none raises it.
The comparison does not isolate word class from the prompts each rule draws or from punctuation, and the results hold for one seed and one corpus.

\subsection{A second target model}\label{sec:attack:leak:llama}

The Corpus A result does not depend on Pythia in the one comparison available.
With a leak of ten prompts and \vf{90} rows, Llama-3.2-1B reaches a word-F1 of \vf{0.632} under the hybrid fit and \vf{0.370} under the supervised fit, against \vf{0.604} and \vf{0.352} for Pythia (Table~\ref{tab:attack:llama}). % src: results/nqllama/v7_2/main/mixed-oct_w5b-s1-prefix.json; results/v7_2/main/mixed-oct_w5b-A-prefix.json
The gain in word-F1 from the matching term is \vf{0.263} for Llama and \vf{0.253} for Pythia, on the \vf{500} Corpus A selection prompts with seed \vf{1}.
The comparison covers one seed, one layer of each model, and two models of about one billion parameters.

\begin{table}[H]
\caption{A second target model reproduces the gain of the matching term on Corpus A.
Word-F1 averaged over three decoders, seed \vf{1}, \vf{500} selection prompts, a leak of \vf{90} rows from ten prompts (whole prompts, the last one in part), hybrid fit with $\alpha = 0.99$, $\tau = d$, $B = 8{,}192$.
The difference is hybrid minus supervised with a \vf{95} percent interval; the mismatched-prompt controls lie between \vf{0.11} and \vf{0.17}.}
\label{tab:attack:llama}
\centering
\begin{tabular}{lccc}
\toprule
Target model & Hybrid fit & Supervised fit & Difference \\
\midrule
Pythia-1B & \vf{0.604} & \vf{0.352} & \vf{$+0.253$} [\vf{0.243}, \vf{0.264}] \\ % src: results/v7_2/main/mixed-oct_w5b-A-prefix.json cells *L90-prefix (0.6042, 0.3516, +0.2526 [0.2426, 0.2640]; controls 0.1063, 0.1433)
Llama-3.2-1B & \vf{0.632} & \vf{0.370} & \vf{$+0.263$} [\vf{0.251}, \vf{0.274}] \\ % src: results/nqllama/v7_2/main/mixed-oct_w5b-s1-prefix.json cells *L90-prefix (0.6324, 0.3696, +0.2628 [0.2514, 0.2742]; controls 0.1544, 0.1680)
\bottomrule
\end{tabular}
\end{table}

Under the hybrid fit, the Corpus A leak of ten prompts yields a word-F1 of \vf{0.604} where the legal128w leak of ten prompts yields \vf{0.304} (Table~\ref{tab:attack:shape}).
The two corpora differ in domain and in prompt length, and the legal leak takes nine rows from each prompt while the Corpus A leak takes whole prompts and part of a last one; the comparison does not separate these causes.

\subsection{A larger decoder and a rewriting model}\label{sec:attack:leak:decoder}

Two further comparisons hold the mapped states of the \vf{90}-prompt leak fixed and change how text is reconstructed from them.
The first contradicts H7: a larger decoder does not raise word-F1 on the same mapped states.
Llama-3.2-3B was attached to the frozen legal128w hub by training one linear projector with the decoding loss of Section~\ref{sec:hub:training}; the hub and the map stayed fixed. % src: oct_w5/decoder-plan.md part 3; results/legal128w/v7_2/main/oct_w5c-projector-llama3b.json training
From the hub outputs of OLMo, the new decoder reaches a word-F1 of \vf{0.983}, against \vf{0.986} for SmolLM, the best hub decoder, so the projector attaches it to the hub. % src: oct_w5c-projector-llama3b.json arms.G2 (0.9833; new_minus_best -0.0023 [-0.0057, 0.0010])
From the mapped states of the hybrid fit, with one row from each of \vf{90} prompts, the new decoder reaches a word-F1 of \vf{0.214} against \vf{0.487} for SmolLM.
The difference is \vf{$-0.273$} (interval \vf{$-0.292$} to \vf{$-0.255$}); its upper bound lies below \vf{0.02}, so the criterion of H7 is not met and the planned prediction is contradicted. % src: same file, arms.b90-a0.99-sq-m0-t0.06485-r1-B8192-L90-random (0.2140; smollm 0.4873; new_minus_best -0.2733 [-0.2918, -0.2551])
The larger decoder scores at the level of Qwen (\vf{0.265}) and OLMo (\vf{0.210}) on those states; SmolLM is the exception among the four. % src: same file, hub_decoders (qwen 0.2648, olmo 0.2099)
The result covers one larger decoder attached by its projector alone, on \vf{500} selection prompts; it does not show that decoder size is irrelevant under a joint fit.

A rewriting model recovers more words from the three reconstructions than their mean word-F1 credits.
An off-the-shelf Qwen3-8B model receives the three decoders' reconstructions of a prompt and writes one text.
On \vf{450} selection prompts, word-F1 rises from \vf{0.340}, the mean of the three reconstructions, to \vf{0.505} after rewriting, a gain of \vf{0.165} (interval \vf{0.155} to \vf{0.176}). % src: results/legal128w/v7_2/main/framework/attacker/attacker-b90-a0.99-sq-m0-t0.06485-r1-B8192-L90-random-s1-stock.json table."hybrid rung B" (0.3400, 0.5052, +0.1652 [0.1547, 0.1759]); prompts_reported 450
Applied to the mismatched-prompt control, the same step gains \vf{0.025} word-F1, so the gain depends on the reconstructions of the prompt itself.
The comparison does not separate words copied from the reconstructions from words that the rewriting model infers from them. % src: same file, table."wrong memory rung B" (+0.0253 [0.0214, 0.0292])
The other \vf{50} selection prompts chose the instruction given to the rewriting model. % src: oct_w5/decoder-plan.md part 2 (selection sample: first 50 prompt ids)
Word-F1 counts overlapping words; the rewriting result does not show that the rewritten text is correct in meaning.
